% \chapter{Практика}

% \section{Задача}
% Разделить сеть 10.4.0.0/16 на сети по 400, 1000, 30 узлов

% \section{Решение}

% Чтобы разбить сеть таким образом, чтобы в каждую подсеть можно было подключить 100 узлов, нужно оставить в узловой части минимум
% 7 бит.

% \subsection{Выделение сети на 1000 узлов}
% Чтобы сеть могла вместить 1000 узлов, нужно оставить в узловой части 10 бит, тогда: \\

% Адрес сети: 10.4.0.0 \\
% Диапазон адресов: 10.4.0.1-10.4.3.254

% \subsection{Выделение сети на 400 узлов}
% Чтобы сеть могла вместить 400 узлов, нужно оставить в узловой части 9 бит, тогда: \\

% Адрес сети: 10.4.4.0 \\
% Диапазон адресов: 10.4.4.1-10.4.5.254

% \subsection{Выделение сети на 30 узлов}
% Чтобы сеть могла вместить 30 узлов, нужно оставить в узловой части 5 бит, тогда: \\

% Адрес сети: 10.4.6.0 \\
% Диапазон адресов: 10.4.6.1-10.4.6.30

% \endinput